% ============================================================
% Cover Letter — BearMamba-3 → IEEE TIM
% Draft: 2026-06-18  (D34, target submission 2026-08-15)
% ============================================================
\documentclass[11pt,a4paper]{article}
\usepackage[margin=1in]{geometry}
\usepackage{amsmath,amssymb}
\usepackage{microtype}
\usepackage{hyperref}
\setlength{\parskip}{6pt}
\setlength{\parindent}{0pt}

\begin{document}

\thispagestyle{empty}

\textbf{Yi Wang}\\
College of Engineering, Lishui University\\
Lishui 323000, Zhejiang, China\\
E-mail: \texttt{jefflsxy@lsu.edu.cn}\\[6pt]
\today

\vspace{12pt}

\textbf{To:} Editor-in-Chief\\
\textit{IEEE Transactions on Instrumentation and Measurement}

\vspace{12pt}

\textbf{Subject:} Submission of original research manuscript ---
``BearMamba-3: Mamba-3 State-Space Bearing Fault Diagnosis with Kinematic
Frequency Inductive Bias and Multi-Sensor Fusion''

\vspace{6pt}

Dear Editor,

We are pleased to submit the manuscript entitled \textbf{``BearMamba-3:
Mamba-3 State-Space Bearing Fault Diagnosis with Kinematic Frequency
Inductive Bias and Multi-Sensor Fusion''} for consideration as a regular
paper in \textit{IEEE Transactions on Instrumentation and Measurement}.
The work is original, has not been previously published, and is not
currently under consideration elsewhere. All authors have approved the
submission.

\textbf{Why IEEE TIM.} Reliable on-line vibration measurement and
condition-monitoring instrumentation lies at the heart of modern
predictive maintenance, and the diagnostic algorithm running on top of
the measurement chain must be both physically interpretable and robust
to changing operating conditions. This work proposes a measurement-stage
building block---an interpretable state-space diagnostic model with
explicit physical frequency alignment---that fits naturally within TIM's
scope of instrumentation, signal-level measurement, and condition-based
diagnostic methods.

\textbf{Core contributions.} The paper makes three contributions:
\begin{enumerate}
\item \textbf{(C1) First adoption of the official Mamba-3 backbone for
bearing diagnostics.} We show that Mamba-3's complex oscillatory states
carry instantaneous rotation frequencies $f_{h,k}(t)$ that are absent in
real-valued SSMs (Mamba-2). This is a structural enabler, not a free
choice of architecture.
\item \textbf{(C2) Kinematic frequency inductive bias $\mathcal{L}_{\text{kin}}$
at the activation level.} The loss aligns Mamba-3 state frequencies with
analytically derived bearing fault harmonics (BPFO, BPFI, BSF, FTF)
computed per sample from the shaft speed. The constraint is undefined
for real-valued SSMs, structurally binding C1 and C2 into a single
integrated design choice. On CWRU, $\mathcal{L}_{\text{kin}}$ reduces
inter-seed variance by $14$--$62\%$ across the EAAI SNR grid.
\item \textbf{(C3) A precision-sample-efficiency trade-off in
multi-sensor fusion.} Dual-sensor fusion yields a statistically
significant $+0.70$~pp gain at $-4$~dB on CWRU
(Wilcoxon $p = 0.031$, $n=5$) and improves inner-race cross-condition
recall by $+31.6$~pp on XJTU-SY (5/5 seeds), while inducing negative
transfer in small-sample leave-one-bearing-out folds with a single
training instance per fault type. This trade-off is reported honestly
with its data-constraint origin clearly stated.
\end{enumerate}

\textbf{Honest scope statement.} We explicitly do \emph{not} claim
within-condition accuracy supremacy on CWRU: a Phase~2b BatchNorm
ablation shows that 1D-CNN with BatchNorm retains a $\sim$10.5~pp edge
over BearMamba-3 at $-8$~dB, and removing BatchNorm closes this gap to
within statistical noise ($p = 0.125$). The paper's distinctive
contribution is in \emph{out-of-distribution} robustness: under
variable-speed cross-condition evaluation on XJTU-SY, BearMamba-3 with
$\mathcal{L}_{\text{kin}}$ achieves macro-F1 of $79.7\%$, $+38.4$~pp
above 1D-CNN, demonstrating that physical inductive bias provides a
form of robustness that statistical normalisation cannot replicate.

\textbf{Concurrent related work.} A concurrent preprint by Li et~al.\
(arXiv:2601.21293, ``Reliability-Calibrated Edge-IoT Early Fault Warning
for Rotating Machinery with a Physics-Guided Tiny-Mamba Transformer'',
submitted to \textit{IEEE Internet of Things Journal}) explores
physics-guided Mamba for edge-IoT deployment and reliability calibration
via extreme value theory. Our work and theirs occupy complementary
niches---we focus on activation-level physical embedding and OOD
robustness in industrial-PC deployments, they focus on sub-1~MB edge
inference with EVT-calibrated false-alarm control. We discuss this
relationship explicitly in Sections~II and~V of the manuscript.

\textbf{Reproducibility.} All experiments are reported with five
independent random seeds, mean $\pm$ standard deviation, Wilcoxon
signed-rank significance, and exhaustive hyperparameter logging.
Negative and boundary-condition results
($\mathcal{L}_{\text{kin}}$ inefficacy on Paderborn discrete-speed and
XJTU-SY end-of-life regimes; LOBO outer-race degradation under single
training bearings) are reported alongside positive findings.
Source code, hyperparameter configurations, and all experimental
artifacts supporting the figures and tables are available for peer
review at:
\par\medskip
\centerline{\url{https://anonymous.4open.science/r/bearmamba3-tim-supplementary-0E81}}
\par\medskip
The package will be deposited to IEEE DataPort upon acceptance with a
citable DOI for permanent archival.

\textbf{Author and conflict-of-interest declaration.} Both authors
contributed substantially to the conception, experimental design,
analysis, and writing. We declare no conflicts of interest. The work was
supported by the Natural Science Foundation of Fujian Province under
Grant 2026J0011040.

\textbf{Suggested reviewers.} We would be glad to suggest the following
experts whose expertise matches the manuscript's scope, while leaving
the final selection to your discretion: (i)~prof.\ \textit{[Suggested
Reviewer 1, name + affiliation + email]}; (ii)~prof.\ \textit{[Suggested
Reviewer 2]}; (iii)~prof.\ \textit{[Suggested Reviewer 3]}. We have no
objection to any reviewer chosen by the editorial board.

We hope the work will be of interest to the readership of \textit{IEEE
Transactions on Instrumentation and Measurement}, and we look forward to
your editorial decision.

\vspace{12pt}

Yours sincerely,\\[12pt]

\textbf{Yi Wang} (corresponding author)\\
College of Engineering, Lishui University\\
Lishui 323000, Zhejiang, China\\
\texttt{jefflsxy@lsu.edu.cn}\\[6pt]
on behalf of co-author \textbf{Yongqing Tang}\\
School of Physics, Mechanics \& Electronics, LongYan University\\
Longyan, Fujian, China\\
\texttt{yq\_tang@lyun.edu.cn}

\end{document}
